\subsection{GRPO background and algorithm overview}

A language model writes a complete response, and a checker or reward model returns one
number for that response. It does not reveal which token was responsible, provide a
corrected response, or expose a derivative that can be backpropagated. The learning
problem is therefore concrete: turn one black-box judgment on a long sampled response
into a parameter update.

The model has only one useful lever. It can change how likely it is to repeat the sampled
decisions. If a response was better than expected, make its token decisions more likely;
if it was worse than expected, make them less likely. This is the intuition behind a
policy gradient. The reward function need not be differentiable because the update
differentiates the model's probability of producing the response, not the checker itself.

Raw reward is not yet a useful direction. A score of $0.8$ may be disappointing for an
easy prompt and excellent for a difficult one. The relevant question is whether a response
was better or worse than the model's usual response to the same prompt. GRPO estimates
that reference point by sampling a group of responses to one prompt and comparing each
response with its siblings. Above-average responses receive positive advantage;
below-average responses receive negative advantage.

Two practical safeguards complete the idea. The responses were generated by a frozen
snapshot of the model, so probability ratios record how the trainable policy has changed
since generation. Clipping stops reused data from encouraging an arbitrarily large move.
A separate reference-policy penalty limits cumulative drift across many new rollout
batches, because a narrow reward can improve while general language behavior degrades.

In one sentence, GRPO samples several responses to the same prompt, turns their relative
scores into positive and negative learning signals, changes the probabilities of their
sampled token decisions, and constrains how far those changes may move the policy. Only
after this causal story is clear is the full formula useful: it compresses these operations
into one scalar minimized by the optimizer.

Let $q$ denote a prompt and let
$o_i=(o_{i,1},\ldots,o_{i,T_i})$ denote sampled response $i$. Using equal trajectory
weighting, the complete loss studied in this section is
\begin{equation}
  \begin{aligned}
  \mathcal{L}_{\mathrm{GRPO}}(\theta)
  =-\frac{1}{G}\sum_{i=1}^{G}\frac{1}{T_i}\sum_{t=1}^{T_i}
  \Big[&\min\!\left(
      \rho_{i,t}(\theta)A_i,
      \operatorname{clip}(\rho_{i,t}(\theta),1-\epsilon,1+\epsilon)A_i
    \right) \\
    &-\beta\big(\exp(d_{i,t})-d_{i,t}-1\big)\Big].
  \end{aligned}
  \label{eq:grpo-full-loss}
\end{equation}
Read the equation as the preceding story in reverse. $A_i$ says whether response $i$ was
better or worse than its siblings. $\rho_{i,t}$ says how the probability of sampled token
$t$ has changed since generation. The minimum applies the clipping guardrail. The term
weighted by $\beta$ is the reference-policy penalty. The two averages combine tokens into
responses and responses into a group; the leading minus sign converts a quantity to
maximize into a loss to minimize.

The three constructed quantities are
\begin{align}
  A_i
  &=\frac{r_i-\mu_q}{\sigma_q+\epsilon_A},
  &\mu_q&=\frac{1}{G}\sum_{j=1}^{G}r_j,
  &\sigma_q^2&=\frac{1}{G}\sum_{j=1}^{G}(r_j-\mu_q)^2,
  \label{eq:grpo-overview-advantage}\\
  \rho_{i,t}(\theta)
  &=\frac{\pi_\theta(o_{i,t}\mid q,o_{i,<t})}
          {\pi_{\mathrm{old}}(o_{i,t}\mid q,o_{i,<t})},
  \label{eq:grpo-overview-ratio}\\
  d_{i,t}
  &=\log\pi_{\mathrm{ref}}(o_{i,t}\mid q,o_{i,<t})
    -\log\pi_\theta(o_{i,t}\mid q,o_{i,<t}).
  \label{eq:grpo-overview-kl-difference}
\end{align}
Here $G$ is the number of sibling responses, $T_i$ is the number of trainable tokens in
response $i$, and $r_i$ is its scalar reward. The small $\epsilon_A$ prevents division by
zero, $\epsilon$ is the clipping width, and $\beta$ controls the reference penalty.
$\pi_{\mathrm{old}}$ is the snapshot that generated the data, while
$\pi_{\mathrm{ref}}$ is a longer-lived anchor.

The next modules rebuild the formula in the order that its needs arise.
Section~\ref{subsec:grpo-group-advantage} explains why sibling comparison creates $A_i$.
Section~\ref{subsec:grpo-gradient-direction} explains how a black-box scalar becomes a
policy gradient. Section~\ref{subsec:grpo-importance-ratio} connects that gradient to the
stored rollout data and to the per-token $\rho_{i,t}A_i$ surrogate.
Section~\ref{subsec:grpo-clipped-objective} motivates clipping and KL before reassembling
Equation~\ref{eq:grpo-full-loss}. The foundational explanation assumes synchronous
generation and a response containing only model-generated tokens; buffers, asynchronous
workers, and agentic traces are introduced afterward. DeepSeek-R1 uses GRPO during its
reinforcement-learning stages \citep{guo2025deepseekr1}.
